\documentclass[11pt]{article}

\usepackage[T1]{fontenc}
\usepackage[utf8]{inputenc}
\usepackage{lmodern}
\usepackage[margin=1in]{geometry}
\usepackage{amsmath,amssymb,amsthm}
\usepackage{booktabs}
\usepackage{tabularx}
\usepackage{array}
\usepackage{enumitem}
\usepackage{forest}
\usepackage[numbers,sort&compress]{natbib}
\usepackage{url}
\usepackage[hidelinks]{hyperref}

\newtheorem{proposition}{Proposition}
\newtheorem{lemma}{Lemma}
\theoremstyle{definition}
\newtheorem{definition}{Definition}

% mana and tap symbols as printed in Oracle text
\newcommand{\ms}[1]{\{#1\}}
\newcommand{\card}[1]{{\mdseries\scshape #1}}
\newcommand{\CR}[1]{CR~#1}
\newcommand{\last}{\mathrm{last}}
\newcommand{\rest}{\mathit{rest}}
\newcommand{\enc}{\mathrm{enc}}
\newcommand{\dec}{\mathrm{dec}}
\newcommand{\mv}{\to}

\forestset{
  board/.style={for tree={grow'=0, parent anchor=east, child anchor=west,
      anchor=west, l sep=5mm, s sep=1.2mm, font=\footnotesize, edge={<-},
      inner sep=1.5pt}},
  leaf/.style={draw, rounded corners=2pt},
  cell/.style={draw, circle, inner sep=0.8pt, font=\scriptsize},
  anch/.style={draw, fill=black!8},
}

\title{Combinators on the Battlefield:\texorpdfstring{\\}{ }BCKW and SKI Reduction in Magic: The Gathering}
\author{N. E. Davis}
% TODO(author): affiliation and email lines, if wanted.
\date{Draft of 21 November 2026}

\begin{document}
\maketitle

\begin{abstract}
We lay out terms of combinatory logic directly as game objects in Magic: The Gathering.
A term is a forest of Aura permanents with ``Enchant permanent'', and application is
attachment: the head of a spine is an Aura, and its arguments hang from a chain of
``cell'' Auras below it. Three primitive game actions, each the printed effect of a real
paper card, reduce such a forest: re-attach an Aura (\card{Simic Guildmage}), copy one Aura
as a token (\card{Flash Photography} or \card{Replicator}), and sacrifice an Aura
(\card{Auratog}). Every rule of B, C, K and W, and of S and I, becomes a fixed sequence of
these actions that realises one normal-order step. W is the only rule of BCKW that copies;
B and C are pure re-attachment; K is one sacrifice, after which the game's state-based
actions delete the discarded argument with no choices. The model is cooperative: one
player chooses every action, the cards permit the reduction moves and do not forbid
others, and only the deletion cascade is automatic. This is not a claim that Magic is
Turing complete in the forced sense of Churchill, Biderman and Herrick. A Python model of
the moves agrees step for step with a reference normal-order reducer on 17 fixed and 500
random terms; it models the moves and selected rule mechanics, not the rules of Magic.
Rules legality rests on our reading of the Comprehensive Rules of 6 November 2026, with
no judge confirmation.
\end{abstract}

\section{Introduction}

Churchill, Biderman and Herrick showed that a two-player game of Magic: The Gathering,
with every move after setup forced, can simulate a universal Turing machine
\citep{cbh2019}. Their tape is a row of creature tokens, the head position is encoded in
power and toughness, and about half the deck exists to remove player choices. Yin and
Churchill later embedded a register machine with a return stack \citep{yin2024}. Both
constructions compile computation into a machine that Magic then simulates.

We ask a different question. Can the terms of a functional calculus be laid out on the
table \emph{as themselves}, so that the board reads as the term and each rewrite rule is a
short sequence of ordinary game actions? We answer it for combinatory logic in the bases
BCKW and SKI.

By \emph{native} we mean four things. (i) Term structure is carried by one relation that
the game already maintains, ``attached to'', and by card names; we use no counters, no
text-changing effects and no position tags. (ii) The redex is located by the shape of the
board, not by counting. (iii) Each primitive action is the printed effect of a real paper
card. (iv) Discarding a subterm is done by the game itself, through state-based actions.
Native is not the same as forced: in our construction a player chooses every move.

\paragraph{Contributions.} Each item carries its status.
\begin{itemize}[nosep,leftmargin=*]
\item The \emph{spine-chain encoding} of terms as a forest of Auras, with a decoding and
  an invariant (Section~\ref{sec:encoding}). \emph{Constructed}; the encoding and
  decoding are inverse by a routine induction.
\item A machine with three primitives (MOVE, COPY, SAC) and a rewrite table for B, C, K,
  W, S and I (Section~\ref{sec:machine}). \emph{Constructed and simulated}; a proposition
  about the abstract moves is given with a proof sketch for K and S and checked for the
  other rows by simulation and an independent hand review.
\item Three card realisations: an 11-card bounded board for $S\,K\,K\,x$, an unbounded
  engine of 13 cards for closed SKI terms, and a pausable 11-card engine built on
  \card{Replicator} (Section~\ref{sec:cards}). \emph{Constructed and simulated at the level of
  the moves}; legality is \emph{argued} from quoted rules and rulings.
\item A measurement of BCKW against SKI on five Church-numeral programs
  (Section~\ref{sec:eval}). \emph{Simulated}; a measurement, not a theorem.
\item Negative results on other designs: the term in the library, the term on the stack,
  sharing, forced reduction and phasing (Section~\ref{sec:fails}). \emph{Argued}, and in
  part \emph{survey results} over card searches that are limited by wording.
\end{itemize}

\paragraph{Non-claims.} We do not claim that the reduction is forced, that Magic is
Turing complete in the sense of \citet{cbh2019} by this route, that the construction is
proved correct as a Magic play, or that any judge has confirmed the rules readings. The
initial board and the engine are assumed to be on the battlefield; we do not construct
them from a deck. The opponent is assumed to pass priority whenever asked.

\section{Background}

\subsection{Combinatory logic}

Terms are built from atoms (combinators and variables) by binary application, written by
juxtaposition and associating to the left \citep{schonfinkel1924,curryfeys1958}. The
rules we use are
\begin{align*}
  B\,x\,y\,z &\to x\,(y\,z) & C\,x\,y\,z &\to x\,z\,y & K\,x\,y &\to x\\
  W\,x\,y &\to x\,y\,y & S\,x\,y\,z &\to x\,z\,(y\,z) & I\,x &\to x .
\end{align*}
Each of $\{S,K\}$ (with $I = S\,K\,K$) and $\{B,C,K,W\}$ is combinatorially complete.
Every term has a unique \emph{spine form} $h\,a_1\cdots a_n$ with $h$ an atom, the
\emph{head}, and $n \ge 0$ arguments. Normal-order (leftmost-outermost) reduction
contracts the head redex when $h$ is a combinator of arity $k \le n$, giving the
contractum followed by $a_{k+1}\cdots a_n$; otherwise it reduces the arguments from left
to right. Normal order finds a normal form whenever one exists. We write $|t|$ for the
number of nodes of $t$ as a binary tree, $2m-1$ for a term with $m$ atoms.

\subsection{The Magic rules we use}

We cite the Comprehensive Rules (CR) of 6 November 2026 \citep{cr2026} and Oracle text and
rulings as served by Scryfall \citep{scryfall}. A separate check compared every card text,
ruling and rule quoted in our construction notes against these sources (24 of 24 cards
and 9 of 9 rulings matched; 46 of 48 rules matched, and the two minor mismatches are
corrected here). Appendix~\ref{app:rules} collects the quoted rules.

\emph{Permanents and Auras.} A permanent is a card or token on the battlefield
(\CR{110.1}). An \emph{Aura} is an enchantment attached to one object; it cannot enchant
itself or more than one object (\CR{303.4d}). All our node cards say ``Enchant
permanent'', so an Aura may enchant another Aura. If an Aura enters ``by any means other
than by resolving as an Aura spell'' and the effect does not say what it enchants, ``that
player chooses what it will enchant as the Aura enters the battlefield'' (\CR{303.4f}); a
ruling on \card{Mythos of Illuna} adds that such an Aura ``doesn't target anything''.

\emph{Tokens and copies.} The creator of a token owns and controls it (\CR{111.2}); a token
outside the battlefield ceases to exist (\CR{111.7}). A copy takes only the copiable values
printed on the object, not ``status, counters, and stickers'' (\CR{707.2}), and the rulings
on \card{Flash Photography} and \card{Replicator} say it does not copy ``any Auras or
Equipment attached to it''. A copy of a spell keeps its targets unless an effect lets its
controller choose new ones (\CR{707.10}, \CR{707.10c}).

\emph{State-based actions, sacrifice, attach.} Whenever a player would get priority, the
game performs all applicable state-based actions as one event and repeats the check until
none applies (\CR{704.3}, \CR{117.5}). The one we rely on is \CR{704.5m}: ``If an Aura is
attached to an illegal object or player, or is not attached to an object or player, that
Aura is put into its owner's graveyard.'' Sacrifice moves a permanent one controls to the
graveyard and is not destruction (\CR{701.21a}). To attach an Aura is ``to take it from
where it currently is and put it onto that object or player'' (\CR{701.3a}).

\emph{The stack} enters only in the copy engines (Section~\ref{sec:cards}). Each new object
goes on top (\CR{405.2}); when all players pass in succession the top object resolves, and
if the stack is empty the step or phase ends (\CR{117.4}, \CR{405.5}).

\subsection{Prior work}

\citet{cbh2019} reduce the halting of Rogozhin's (2,18) universal Turing machine to the
outcome of a forced two-player game, so non-halting appears as a mandatory loop, a draw
under \CR{104.4b}. Churchill's earlier machines (2010--2014) needed ``may'' choices
\citep{churchillsite}. \citet{yin2024} embed a register machine with 12 registers,
addressed memory and a return stack; the program is the order of lands in a library, and
the return stack is a chain of tokens linked by phasing. \citet{biderman2020} shows that
deciding mate-in-$n$ is $\Delta^0_n$-hard and that optimal play is non-arithmetic in general. \citet{chatterjee2016} show
that deciding the legality of a single step is coNP-complete in general (we know this
paper from a search result only). For other card games, \citet{nicolosi2026} prove an
undecidability result for strategies in Yu-Gi-Oh!, and \citet{dalaklis2026} claim Turing
completeness for Riftbound.

Our literature search was six web-search queries (for lambda calculus, combinatory logic,
SKI, term rewriting, stack machines, Church numerals and related terms in Magic) through
a general search engine, plus the searches behind the sources above. It found no
embedding of lambda calculus, combinatory logic or term rewriting in Magic. The search did
not reach forum or chat archives (Reddit, MTG Salvation, Discord, Stack Exchange)
directly, so an absence here does not exclude forum posts. We read the CBH paper only
through a summarised HTML rendering.

\section{The encoding}\label{sec:encoding}

A \emph{board} is a finite set of Aura permanents, the \emph{nodes}, together with one
non-Aura permanent $A$, the \emph{anchor}. Each node $n$ has a \emph{label} $\lambda(n)$,
which is its card name, and a \emph{host} $\eta(n)$, the permanent it is attached to. The
label is either an atom or the special label \emph{cell}; nodes with atom labels are
\emph{leaves}. Every label card is a distinct Aura with ``Enchant permanent''.

\begin{definition}[spine-chain encoding]\label{def:enc}
For a term $t$ and a permanent $p$, $\enc(t,p)$ adds fresh nodes to the board as follows.
Write $t = h\,a_1\cdots a_n$ in spine form. Add a leaf $\ell$ with $\lambda(\ell)=h$ and
$\eta(\ell) = p$. If $n \ge 1$, add cells $c_1,\dots,c_n$ with $\eta(c_1)=\ell$ and
$\eta(c_{i+1}) = c_i$, and for each $i$ perform $\enc(a_i, c_i)$. The board of a closed
or open term $t$ is $\enc(t, A)$.
\end{definition}

So the head is the one leaf on the anchor, and argument $i$ hangs from the cell at depth
$i$ of the head's chain (Figure~\ref{fig:enc}). An application node costs one cell and an
atom one leaf, so $t$ takes $|t|$ Auras.

\begin{definition}[well-formed board; decoding]\label{def:wf}
A board is \emph{well formed} if (i) the host relation has no cycle and every node reaches
$A$; (ii) every leaf enchants a cell or $A$, and every cell enchants a leaf or a cell;
(iii) $A$ has exactly one child, a leaf; each node has at most one cell child; each cell
has exactly one leaf child; and no leaf has a leaf child. For a leaf $\ell$ with cell
chain $c_1,\dots,c_k$ (each $c_{i+1}$ the cell child of $c_i$), let
$\dec(\ell) = \lambda(\ell)\,\dec(\mathrm{leaf}(c_1))\cdots\dec(\mathrm{leaf}(c_k))$,
and let the board decode to $\dec$ of the leaf on $A$.
\end{definition}

\begin{lemma}
$\enc(t,A)$ is well formed and decodes to $t$. Conversely, a well-formed board that decodes to
$t$ equals $\enc(t,A)$ up to the identity of the objects.
\end{lemma}
\noindent The proof is a routine induction on $t$ and on the height of the board.

We use two derived notions. For a leaf $r$, $\last(r)$ is the last cell of its chain, or
$r$ itself if it has none. The \emph{rest} of a redex $h\,a_1\cdots a_k\,a_{k+1}\cdots$ is
the cell $c_{k+1}$ with everything below it; it is absent when the redex has exactly $k$
arguments. In the figures an arrow points from an Aura to the permanent it enchants.

\begin{figure}[t]
\centering
\begin{forest} board
[$A$,anch [$h$,leaf [$c_1$,cell [$a_1$,leaf] [$c_2$,cell [$a_2$,leaf] [$c_3$,cell [$a_3$,leaf]]]]]]
\end{forest}
\hfill
\begin{forest} board
[\card{Island},anch [$S$,leaf [$c_1$,cell [$K$,leaf] [$c_2$,cell [$K$,leaf] [$c_3$,cell [$x$,leaf]]]]]]
\end{forest}
\caption{Left: the encoding of $h\,a_1\,a_2\,a_3$ when each $a_i$ is an atom; a compound
$a_i$ hangs its own leaf and chain from $c_i$. Right: $S\,K\,K\,x$. Boxes are leaves, circles
are cells, the shaded node is the anchor. Within a node's children the leaf is drawn
first.}\label{fig:enc}
\end{figure}

\section{The machine}\label{sec:machine}

\subsection{Three primitives}

\begin{description}[leftmargin=1em,nosep]
\item[MOVE$(n \mv p)$] re-attaches node $n$ to permanent $p$. Card: \card{Simic Guildmage},
  ``\ms{1}\ms{U}: Attach target Aura attached to a permanent to another permanent with the
  same controller.'' Its ruling of 2006-05-01 says the new host is chosen on resolution, ``can't be
  the permanent the Aura is already attached to, it must be controlled by the player who
  controls the permanent the Aura is attached to, and it must be able to be enchanted by
  the Aura.'' Rules: \CR{701.3a}.
\item[COPY$(n \mv p)$] creates a token copy of node $n$ attached to $p$. Card:
  \card{Flash Photography}, ``Create a token that's a copy of target permanent'', or
  \card{Replicator}, ``\ms{2}, \ms{T}: Create a tapped token that's a copy of target token you
  control.'' The token has $n$'s label (\CR{707.2}, \CR{111.3}) and no attachments (rulings);
  its host is chosen as it enters (\CR{303.4f}).
\item[SAC$(n)$] sacrifices node $n$. Card: \card{Auratog}, ``Sacrifice an enchantment: This
  creature gets +2/+2 until end of turn.'' Sacrifice is a cost, so it needs no target.
  Rules: \CR{701.21a}, then \CR{704.3} and \CR{704.5m}.
\end{description}

When a node moves, the nodes attached to it stay attached and travel with it. No rule says
so in positive terms. The argument is from absence: attaching changes only what the named
Aura is put onto (\CR{701.3a}); it is not a zone change, so the Aura remains the same object
(contrast \CR{400.7}); and none of the listed ways for an attachment to end
(\CR{303.4d}, \CR{701.3d}, \CR{704.5m}, \CR{704.5p}) applies to an Aura whose host is still on
the battlefield.

\subsection{The rewrite table}

Let $h$ be the head leaf of the normal-order redex and $P = \eta(h)$, which is $A$ or, when
normal order works inside an argument, the cell that holds that argument. Let
$c_1,\dots,c_k$ be the redex cells, $r_j = \mathrm{leaf}(c_j)$ the root leaf of argument
$j$, and $\rest = c_{k+1}$. Table~\ref{tab:rewrite} gives each rule. The bracketed move is
made only when $\rest$ is present. Each sequence is followed by its copies (top-down, node
by node) and ends with SAC $h$. The table is the one in our construction notes, generated
from the simulator's \texttt{rewrite()} function; only the row order is changed.

\begin{table}[t]
\centering\small
\begin{tabularx}{\textwidth}{@{}l>{\raggedright\arraybackslash}Xcc>{\raggedright\arraybackslash}p{3.0cm}@{}}
\toprule
rule & MOVE sequence (then COPY, then SAC $h$) & MOVEs & COPYs & dies in the cascade\\
\midrule
$B\,x\,y\,z \to x\,(y\,z)$ & $r_1\mv P$; $c_2\mv\last(r_1)$; $c_3\mv\last(r_2)$; $[\rest\mv c_2]$ & 3 (+1) & 0 & $c_1$ (empty)\\
$C\,x\,y\,z \to x\,z\,y$ & $r_1\mv P$; $c_3\mv\last(r_1)$; $c_2\mv c_3$; $[\rest\mv c_2]$ & 3 (+1) & 0 & $c_1$ (empty)\\
$K\,x\,y \to x$ & $r_1\mv P$; $[\rest\mv\last(r_1)]$ & 1 (+1) & 0 & $c_1$, $c_2$ and the whole subtree of $y$: $\mathrm{depth}(y)+2$ passes\\
$W\,x\,y \to x\,y\,y$ & $r_1\mv P$; $c_2\mv\last(r_1)$; $c_1\mv c_2$; $[\rest\mv c_1]$; COPY subtree$(r_2)$ onto $c_1$ & 3 (+1) & $|y|$ & nothing\\
\midrule
$S\,x\,y\,z \to x\,z\,(y\,z)$ & $r_1\mv P$; $c_3\mv\last(r_1)$; $c_2\mv c_3$; $[\rest\mv c_2]$; $c_1\mv\last(r_2)$; COPY subtree$(r_3)$ onto $c_1$ & 4 (+1) & $|z|$ & nothing\\
$I\,x \to x$ & $r_1\mv P$; $[\rest\mv\last(r_1)]$ & 1 (+1) & 0 & $c_1$ (empty)\\
\bottomrule
\end{tabularx}
\caption{Rewrite table. ``COPY subtree$(r)$ onto $c$'' copies the leaf $r$ onto $c$, then
for each cell of $r$'s chain copies the cell onto the previous new node and recursively
copies that cell's leaf onto the new cell. Every token enters attached to a node that
already exists. In the K row, $\mathrm{depth}(y)$ counts the Auras on the longest attachment
path of $y$'s encoding.}\label{tab:rewrite}
\end{table}

An earlier, shortened version of this table dropped the bracketed $\rest$ moves. It was
wrong for every redex with trailing arguments: $K\,x\,y\,r$ lost $r$. The rest moves are
part of each rule.

Within BCKW the division of labour is clean. B and C only re-attach. K is one move and one
sacrifice, and the game deletes $y$. W is the only rule that copies. In SKI, S both
rearranges and copies, and K does the deletion.

\subsection{Correctness of the abstract moves}

We state correctness for the abstract machine under three assumptions.
(M1) MOVE$(n\mv p)$ changes $\eta(n)$ to $p$ and nothing else.
(M2) COPY$(n\mv p)$ adds a fresh node with label $\lambda(n)$, host $p$ and no children.
(M3) SAC$(n)$ removes $n$; then, repeatedly, every node whose host has been removed is
removed, one level per pass, until none is left.
(M1)--(M3) are our reading of \CR{701.3a}, \CR{707.2} with the rulings, and \CR{704.3} with
\CR{704.5m}; they are not theorems about Magic.

\begin{proposition}\label{prop:rows}
Let a well-formed board decode to $t$, and let the normal-order redex of $t$ be an
instance of one of the six rules, with head leaf $h$ on host $P$. Apply that rule's row of
Table~\ref{tab:rewrite} under (M1)--(M3). Then:
(a) in every state before SAC $h$, every node has a host on the board and the host relation
has no cycle, so no node is removed before SAC $h$;
(b) after SAC $h$ and its cascade the board is well formed and decodes to the result
$t'$ of contracting that redex in $t$;
(c) the cascade removes exactly the nodes in the last column of the table.
\end{proposition}

\begin{proof}[Proof sketch]
Two facts are used throughout. A MOVE of $n$ to $p$ creates a cycle exactly when $p$ lies
in the subtree of $n$; and a MOVE leaves every node attached, because it only changes one
host to another node that is present. The rewrite touches only $h$, its chain and the
subtrees hanging from it; $P$'s other child (the next cell of an outer chain, if $P$ is a
cell) is untouched, so the rest of the board decodes as before.

\emph{K.} Before: $\eta(h)=P$, $\eta(c_1)=h$, $\eta(r_1)=c_1$, $\eta(c_2)=c_1$,
$\eta(r_2)=c_2$, and, if present, $\eta(c_3)=c_2$. Move $r_1$ to $P$: $P$ is an ancestor of
$r_1$, not a descendant, so no cycle; $r_1$ carries all of $x$. If $\rest=c_3$ is present,
move it to $\last(r_1)$: $\last(r_1)$ lies in the subtree of $r_1$, which is disjoint from
the subtree of $c_3$, so no cycle; $c_3$ carries the arguments $a_3,\dots$ and they now
extend $x$'s chain. All nodes still have hosts, which proves (a). Now $P$ has leaf child
$r_1$, whose chain is $x$'s own cells followed by $c_3,\dots$, so $r_1$ decodes to
$x\,a_3\cdots a_n$, the contractum with the rest appended. Sacrifice $h$. Its only child
was $c_1$, so pass 1 removes $c_1$; pass 2 removes $c_2$, now hostless; the later passes
remove the encoding of $y$ level by level. That is $\mathrm{depth}(y)+2$ passes and exactly
$c_1$, $c_2$ and $y$, which proves (c). What remains is well formed, which proves (b).

\emph{S.} Before: $h$ on $P$; chain $c_1,c_2,c_3$ with $r_1,r_2,r_3$ (encoding $x,y,z$),
and possibly $c_4=\rest$ on $c_3$. The moves, in order:
(1) $r_1\mv P$, as in K.
(2) $c_3\mv\last(r_1)$: $c_3$ carries $r_3$ and $c_4$; $\last(r_1)$ is in $x$'s subtree,
disjoint from $c_3$'s; no cycle.
(3) $c_2\mv c_3$: since step (2), $c_2$'s only child is $r_2$, so $c_3$ is not below $c_2$;
no cycle. Now $c_3$ has two cell children, $c_2$ and $c_4$; every node still has a host.
(4) $[c_4\mv c_2]$: $c_2$ is not below $c_4$; no cycle. The chain under $r_1$ is now
$x$'s cells, then $c_3$ (holding $z$), then $c_2$ (holding $y$), then $c_4,\dots$.
(5) $c_1\mv\last(r_2)$: $c_1$ has no children left, so no cycle; it extends $y$'s chain.
(6) COPY the subtree of $r_3$ onto $c_1$, top-down; each token enters on a node that exists
and has no children, so no cycle.
This proves (a). Now $c_2$ decodes to $y\,z'$ with $z'$ a copy of $z$, and the board
decodes to $x\,z\,(y\,z)\,a_4\cdots a_n$, which proves (b). $h$ has no children when it is
sacrificed, so the cascade removes nothing, which proves (c).

\emph{I} is the K case without $c_2$ and $y$: the cascade removes the empty $c_1$.
\emph{B, C, W.} The check is the same bookkeeping; in each, every moved cell has been
emptied, or carries only what should travel with it, before it moves. We have not written
it out here. It was confirmed in lock step by the simulator (Section~\ref{sec:eval}) and by
an independent hand review of the general table, which also checked trailing arguments
and nested redexes.
\end{proof}

The proposition says nothing about Magic legality beyond (M1)--(M3). In particular the
intermediate states break invariant (iii) of Definition~\ref{def:wf} (in S, $c_3$
briefly has two cell children), but never (i), which is what keeps state-based actions
quiet until the sacrifice.

\paragraph{Where label effects land.} Leaves move only to $P$, a cell or the anchor; cells
move only to $\last(r)$, a cell or a leaf, or to another cell; copies of leaves land on
cells and copies of cells on leaves or cells. So at every point a leaf enchants a cell or
the anchor and a cell enchants a leaf or a cell. Section~\ref{sec:cards} uses this to show
that the label cards' own effects do nothing.

\subsection{Cost per rule}

Table~\ref{tab:cost} gives the cost of one head step, measured by the simulator on copied
arguments of $n = 1, 2, 4, 8$ atoms. MOVEs are constant per rule. COPYs are $2n-1$ for an
$n$-atom copied argument. Each step has one SAC. Mana depends on the engine (Section~\ref{sec:cards}):
$2\cdot\mathrm{MOVE} + 6\cdot\mathrm{COPY}$ for the \card{Flash Photography} engine and
$2\cdot(\mathrm{MOVE}+\mathrm{COPY})$ for the \card{Replicator} engine. Re-attachment costs
the same whatever the size of the subterm moved.

\begin{table}[t]
\centering\small
\begin{tabular}{@{}lcccc@{}}
\toprule
head step & MOVEs & COPYs & mana, FP engine & mana, Replicator engine\\
\midrule
$I\,x$, $K\,x\,y$ & 1 & 0 & 2 & 2\\
$B\,x\,y\,z$, $C\,x\,y\,z$ & 3 & 0 & 6 & 6\\
$W\,x\,y$, $y$ of $n$ atoms & 3 & $2n-1$ & $6+6(2n-1)$ & $2(3+2n-1)$\\
$S\,x\,y\,z$, $z$ of $n$ atoms & 4 & $2n-1$ & $8+6(2n-1)$ & $2(4+2n-1)$\\
\midrule
with a further argument & $+1$ & & $+2$ & $+2$\\
\bottomrule
\end{tabular}
\caption{Cost of one head step (from the per-rule tables in \texttt{sim/L6/output.txt} and
\texttt{sim/L9/output.txt}, which list $n=1,2,4,8$; for example $S\,x\,y\,z$ with $n=8$:
4 moves, 15 copies, 98 and 38 mana). The K cascade takes 3, 5, 8 and 12 passes on the
measured arguments $y$ of 1, 2, 4 and 8 atoms.}\label{tab:cost}
\end{table}

\section{Card realisations}\label{sec:cards}

All three realisations make the same assumptions. One player, $P$, controls every permanent
involved: the engine, the anchor and every node. The opponent passes priority whenever
$P$ passes. The engine and the initial board are on the battlefield at the start; we do not
construct them from a deck. Appendix~\ref{app:cards} gives the full Oracle text of every
card.

\paragraph{Single controller.} The \card{Simic Guildmage} ruling requires the new host to be
controlled by the player who controls the old host. Every node is a host at some time, so
all nodes and the anchor need one controller. With $P$ controlling everything the clause
always holds. The variable labels change control (``You control enchanted permanent''),
which is idempotent here: $P$ already controls the host.

\paragraph{Labels and why they are inert.} Table~\ref{tab:labels} lists the label cards.
Each is a non-legendary Aura with ``Enchant permanent'', and each has an effect of its own.
By the invariant above, a leaf label's effect falls only on a cell or on the anchor, and
the cell label's effect falls only on nodes. The table gives, for each, why the effect does
nothing there.

\begin{table}[t]
\centering\small
\begin{tabularx}{\textwidth}{@{}l>{\raggedright\arraybackslash}p{3.1cm}>{\raggedright\arraybackslash}X@{}}
\toprule
role & card & why its effect is inert where it lands\\
\midrule
cell & \card{Indestructibility} & nothing is ever destroyed; sacrifice and \CR{704.5m} are not destruction (\CR{701.21a})\\
S & \card{Bound in Gold} & ``activated abilities can't be activated unless they're mana abilities'': cells have no activated abilities and the Island's is a mana ability\\
K & \card{Intercessor's Arrest} & the same\\
I & \card{Curator's Ward} & hexproof stops only opponents (\CR{702.11b}); cells and basic lands are not historic (\CR{700.6}), so under the intervening-if rule (\CR{603.4}) the draw trigger never fires\\
B & \card{Essence Leak} & lands only on white cells or the colourless anchor; ruling: ``It can enchant a permanent that is not red or green, but it doesn't do anything in that case.''\\
C & \card{Faith's Fetters} & ETB life gain each time a C node is copied; harmless\\
W & \card{Elemental Resonance} & triggers only at the start of the first main phase and adds mana\\
variables & \card{Confiscate}, \card{Lay Claim}, \card{Volition Reins}, \card{Take Possession}, \card{Psychic Overload}, \card{Dream Leash} & control change is idempotent under one controller; tapping or untapping a cell is harmless; split second and \card{Dream Leash}'s casting restriction matter only when the card is cast\\
\bottomrule
\end{tabularx}
\caption{Label Auras for the \card{Flash Photography} engine. A closed SKI deck needs the first
four rows; a BCKW deck can avoid all triggers with C~=~\card{Bound in Gold} and
W~=~\card{Curator's Ward}. Only the mixed seven-label deck uses \card{Faith's Fetters} and
\card{Elemental Resonance}.}\label{tab:labels}
\end{table}

The label search also rejected cards that would damage the structure. For example,
\card{Song of the Dryads} turns a node into a land, which as an attached non-Aura becomes
unattached (\CR{704.5p}); the Runes draw a card on entering, which would deck $P$ under
unbounded copying; and Auras with ``Enchant nonland permanent'' or ``Enchant enchantment''
cannot sit on the anchor land, where any label may end up as the head. (These rejected
cards were not part of the source check.)

\subsection{\texorpdfstring{A bounded board: $S\,K\,K\,x$ in 11 cards}{A bounded board: S K K x in 11 cards}}

The cards are an Island (anchor), \card{Bound in Gold} (S), two \card{Intercessor's Arrest}
(K), three \card{Indestructibility} (cells), \card{Confiscate} ($x$), \card{Simic Guildmage},
\card{Auratog} and \card{Flash Photography}. The seven Auras are cast onto each other to
form Figure~\ref{fig:enc} (right); the run takes five moves, one copy and two sacrifices
(Section~\ref{sec:skk}). By the printed costs that is five activations at \ms{1}\ms{U} and one
cast of \card{Flash Photography} at \ms{2}\ms{U}\ms{U}, 14 mana in all. The mana and the
initial layout are assumed. Each copied node uses one cast of a copy spell, and
\card{Flash Photography} gives two (cast and flashback), so this board is bounded.

\subsection{The unbounded engine}

To copy without bound, the second construction parks one \card{Flash Photography} (FP) on
the stack and copies it repeatedly.
\begin{itemize}[nosep,leftmargin=*]
\item \emph{Open a session.} In $P$'s main phase, with the stack empty, cast FP targeting a
  spare Island $T$. The stack is $[\mathrm{FP}(T)]$.
\item \emph{COPY$(n\mv p)$.} Cast \card{Reiterate} with buyback (``Copy target instant or
  sorcery spell. You may choose new targets for the copy.'') targeting FP. When it resolves,
  an FP copy with new target $n$ goes on the stack (\CR{707.10c}) and \card{Reiterate}
  returns to hand (\CR{702.27a}). When the copy resolves, a token copy of $n$ enters
  attached to $p$, chosen under \CR{303.4f}. Reiterate needs FP on the stack both when cast
  and when it resolves (\CR{608.2b}), so FP must stay on the stack for the whole session.
\item \emph{Mana.} \card{Isochron Scepter} imprinted with \card{Dramatic Reversal}, and
  \card{Gilded Lotus}. One loop taps the Lotus for three mana of one colour, spends \ms{2}
  on the Scepter, and casts the free copy of \card{Dramatic Reversal}, which untaps both.
  The net gain is one mana of a colour chosen afresh each loop, and the stack returns to
  what it was. Reiterate with buyback costs 6 mana including \ms{R}\ms{R}, so a COPY
  costs 6 loops; a MOVE costs 2; a SAC costs nothing. Mana empties only at the end of a step
  or phase (\CR{106.4}), and the whole computation fits inside one main phase.
\item \emph{Close the session.} With only FP on the stack, all pass; FP resolves and copies
  the Island $T$, which is harmless. A stray pass with FP alone on the stack ends the
  session early; flashback allows one restart.
\end{itemize}
The engine is \card{Flash Photography}, \card{Reiterate}, \card{Isochron Scepter},
\card{Dramatic Reversal}, \card{Gilded Lotus}, \card{Simic Guildmage}, \card{Auratog} and two
Islands. With four labels that is 13 cards for closed SKI terms; BCKW alone takes 14, and
the mixed six-combinator basis 16. Open terms add one variable label each, up to six. In
this engine the term itself is built from physical template Auras by COPYs (62 mana for
$S\,K\,K\,x$ in simulation). Every action is optional, so a reduction that never ends is
not a mandatory loop and not a draw (\CR{104.4b}); ``halting'' means only that $P$ sees no
redex and stops.

\subsection{The Replicator engine}

\card{Replicator} (\ms{2}, \ms{T}: ``Create a tapped token that's a copy of target token you
control'') replaces the parked spell. \card{Dramatic Reversal} untaps it with everything else
(``Untap all nonland permanents you control''), so each COPY costs \ms{2}, or two loops,
and one stack object. Every primitive is an activated ability that leaves the stack as it
found it, so the run can stop at any primitive boundary, let the step or turn end, and
resume. A pause loses at most the unspent mana in the pool (\CR{106.4}).

The engine is \card{Replicator}, \card{Lithoform Engine}, \card{Isochron Scepter},
\card{Dramatic Reversal}, \card{Gilded Lotus}, \card{Simic Guildmage} and \card{Auratog}. No land
is needed: \card{Replicator}'s Food token hosts the label cards, and \card{Lithoform Engine}
is the anchor. With four labels that is 11 cards for closed SKI terms, and 14 for the
mixed basis. Open terms are limited to five variables: split second on \card{Take
Possession} forbids the activation that would seed it. Because \card{Bound in Gold} and
\card{Intercessor's Arrest} on the anchor switch off \card{Lithoform Engine}'s activated
abilities, every label must be seeded before the term is built.

\paragraph{Seeding caveat.} \card{Replicator} copies only tokens, and the label cards are
cards. So each label needs one card-to-token step. We cast the label card and copy the Aura
spell with \card{Lithoform Engine} (``\ms{4}, \ms{T}: Copy target permanent spell you
control. (The copy becomes a token.)''); the copy becomes a token permanent as it resolves
(\CR{707.10f}). That the token is then a valid seed rests on our reading of two rulings.
\card{Replicator}'s first ruling says ``If the original token is copying something else,
then the new token enters as whatever that token copied''; a \card{Lithoform Engine} ruling
says ``The token that a resolving copy of a spell becomes isn't said to have been
`created'.'' We read the first as covering this case; no judge has confirmed it. The
fallback is \card{Flash Photography} plus \card{Mythos of Illuna}, at one more card.
\card{Replicator} was released on 13 November 2026, and further rulings may follow.

A variant without \card{Simic Guildmage} (engine R) rebuilds each reduct by copies and
sacrifices the old head, so the cascade removes the old subterm. It takes 10 cards for
closed SKI and much more mana (Section~\ref{sec:eval}).

\section{Worked example}\label{sec:example}

We follow trace~2 of \texttt{sim/L6/output.txt}: $S\,x\,y\,(K\,a\,b) \to x\,a\,(y\,a)$.
Object numbers are the simulator's; variables use the labels \card{Volition Reins} ($x$),
\card{Take Possession} ($y$), \card{Confiscate} ($a$) and \card{Lay Claim} ($b$).
Figure~\ref{fig:example} shows the board before, after step 1 and at the end.

\emph{Step 1} (head S\#14 on the anchor). Four moves:
x\#16 from c15 to the anchor ($x$ becomes the head);
c19 from c17 to x\#16 (the cell holding $z$);
c17 from c15 to c19 (the cell holding $y$);
c15 from S\#14 to y\#18 (the freed cell, the argument of $y$).
Then five copies rebuild $z = K\,a\,b$ top-down on c15: K\#25 on c15, c26 on K\#25,
a\#27 on c26, c28 on c26, b\#29 on c28. S\#14 is sacrificed; nothing hangs from it, so no
state-based action fires. The board reads $x\,(K\,a\,b)\,(y\,(K\,a\,b))$.

\emph{Step 2} (head K\#20 on c19). a\#22 moves from c21 to c19; K\#20 is sacrificed. Three
cascade passes remove c21, then c23, then b\#24. The board reads $x\,a\,(y\,(K\,a\,b))$.

\emph{Step 3} (head K\#25 on c15, inside the argument of $y$). a\#27 moves from c26 to
c15; K\#25 is sacrificed; three passes remove c26, c28 and b\#29. The board reads
$x\,a\,(y\,a)$, the normal form.

Totals from the output file: 3 steps, 6 moves, 5 copies, 3 sacrifices, 6 cascade passes,
42 mana, 104 stack resolutions, and at most 15 nodes on the board. Steps 2 and 3 show
normal order moving inside a term: the head's host $P$ is a cell, not the anchor.

\begin{figure}[t]
\centering
\begin{forest} board
[{Island\#5},anch [S\#14,leaf [15,cell [x\#16,leaf] [17,cell [y\#18,leaf] [19,cell [K\#20,leaf [21,cell [a\#22,leaf] [23,cell [b\#24,leaf]]]]]]]]]
\end{forest}\\[2mm]
\begin{forest} board
[{Island\#5},anch [x\#16,leaf [19,cell [K\#20,leaf [21,cell [a\#22,leaf] [23,cell [b\#24,leaf]]]] [17,cell [y\#18,leaf [15,cell [K\#25,leaf [26,cell [a\#27,leaf] [28,cell [b\#29,leaf]]]]]]]]]]
\end{forest}\\[2mm]
\begin{forest} board
[{Island\#5},anch [x\#16,leaf [19,cell [a\#22,leaf] [17,cell [y\#18,leaf [15,cell [a\#27,leaf]]]]]]]
\end{forest}
\caption{$S\,x\,y\,(K\,a\,b)$: before (top), after step 1, which reads
$x\,(K\,a\,b)\,(y\,(K\,a\,b))$ (middle), and at the end, $x\,a\,(y\,a)$ (bottom). Circles
are cells, numbered by object id. Nodes 25--29 are the tokens made by the five copies of
step 1.}\label{fig:example}
\end{figure}

\subsection{\texorpdfstring{$S\,K\,K\,x$}{S K K x}}\label{sec:skk}

Trace~1 of the same file runs $S\,K\,K\,x$ from Figure~\ref{fig:enc} (right). Step 1 moves
K\#13 onto the Island, c16 onto that K, c14 onto c16, and c12 onto the second K (K\#15); it
copies x\#17 to a new x\#18 on c12 and sacrifices S\#11. The board reads $K\,x\,(K\,x)$.
Step 2 moves x\#17 onto the Island and sacrifices K\#13; five successive cascade passes
remove c16, c14, K\#15, c12 and x\#18. Totals: 2 steps, 5 moves, 1 copy, 2 sacrifices,
5 passes, 16 mana and 42 stack resolutions with the \card{Flash Photography} engine, and
12 mana and 32 resolutions with the \card{Replicator} engine.

\section{Evaluation}\label{sec:eval}

Two Python programs, included in \texttt{sim/}, model the moves. They model the abstract
moves and a selection of rule mechanics; they do not model the rules of Magic, and they do
not show that the mechanics they check are all the relevant ones.

\paragraph{What is modelled and checked.} \texttt{l6sim.py} models the battlefield, the
stack (FP, Reiterate and its copies, the Guildmage, Auratog and Scepter abilities, Dramatic
Reversal copies, ETB triggers), the mana pool by colour, and the state-based cascade. It
checks a set of rules conditions (Reiterate needs FP on the stack and Reiterate in hand; an
FP copy's target must be on the battlefield when it resolves; a host must exist under
\CR{303.4f}; the Guildmage's new host must differ from the current one; sorcery-speed casts
need an empty stack; mana must be paid) and protocol conditions (acyclicity after every
primitive; the encoding invariants and the leaf half of the label invariant between
steps). After every step it decodes the board and asserts equality with a reference
leftmost-outermost reducer, and it compares each normal form with a second normaliser that
shares the contraction code and differs only in strategy. \texttt{l9sim.py} reuses the
reference reducer and rule sequences, replaces COPY by \card{Replicator} (token-only target,
untapped source, tapped result), adds \card{Lithoform Engine} seeding, engine R and pauses,
and asserts both halves of the leaf/cell invariant after every primitive.

\paragraph{What is not checked.} The \card{Simic Guildmage} same-controller clause, the
legality of attachment cycles, and trigger timing; the opponent (always passes); the setup of
the engine; and every rules claim in this paper, which rests on quoted text.

\paragraph{Results.} From \texttt{sim/L6/output.txt}: 17 fixed terms all agree in lock step
(among them $S\,K\,K\,x$, $S\,x\,y\,(K\,a\,b)$, Church numerals, $K\,(I\,a)$, which reduces
under an argument, and $K\,a\,(S\,I\,I\,(S\,I\,I))$ and $K\,a\,(W\,W\,W)$, which discard a
divergent argument with one move, one sacrifice and a 10- or 6-pass cascade); three
non-normalising terms ($S\,I\,I\,(S\,I\,I)$, $W\,W\,W$ and $Y\,f$) agree on every one of 40
steps; and 500 random terms of 1 to 24 atoms over $S,K,I,B,C,W,a,b,c$ (seed 20261121) all
agree in lock step, 496 reaching normal form and 4 stopping at the 80-step bound, for a
total of 1796 steps, 4918 moves, 5693 copies and 1796 sacrifices. From
\texttt{sim/L9/output.txt}: the same fixed, non-normalising and random terms agree in both
engines, the random set also with random pauses (1816 pauses, 898 of them passing a
turn, in the engine-G run), and engine G's move, copy and sacrifice counts equal the
\card{Flash Photography} engine's on every fixed term. Both programs end with \texttt{ALL
OK}; we re-ran both from \texttt{sim/} and obtained identical output. The independent
review of the rewrite table reported a further 3000 random terms in lock step and
mutation tests that caught wrong move sequences; that run is not archived here.

\paragraph{Basis comparison.} Five lambda programs were compiled three ways and run to
normal form: SKI by bracket abstraction with K and eta rules; BCKW, which uses B and C when
the variable occurs on one side only, and $W\,(B\,(C\,M')\,N')$ when it occurs on both;
and SKIBC in Turner's style \citep{turner1979}. Table~\ref{tab:basis} gives the results.
On these five programs BCKW uses 25--67\% less mana than SKI with the \card{Flash
Photography} engine, mainly because it copies less, and SKIBC uses less still. This
compares one SKI compilation with B/C-aware compilations on five programs; it is a
measurement, not a statement about the bases.

\begin{table}[t]
\centering\small
\begin{tabular}{@{}lccc@{}}
\toprule
program & SKI & BCKW & SKIBC\\
\midrule
succ 2 $p$ $q$   & 112 / \phantom{0}88 \ (38/6)\phantom{00}   & \phantom{0}84 / \phantom{0}76 \ (36/2)\phantom{0} & \phantom{0}48 / \phantom{0}40 \ (18/2)\phantom{0}\\
plus 2 2 $p$ $q$ & 434 / 250 \ (79/46)\phantom{0}  & 156 / 144 \ (69/3)\phantom{0} & \phantom{0}90 / \phantom{0}78 \ (36/3)\phantom{0}\\
mult 2 2 $p$ $q$ & 514 / 270 \ (74/61)\phantom{0}  & 188 / 136 \ (55/13) & 122 / \phantom{0}86 \ (34/9)\phantom{0}\\
2 2 $p$ $q$      & 634 / 322 \ (83/78)\phantom{0}  & 286 / 190 \ (71/24) & 182 / 118 \ (43/16)\\
mult 3 2 $p$ $q$ & 986 / 478 \ (112/127)          & 324 / 224 \ (87/25) & 206 / 138 \ (52/17)\\
\bottomrule
\end{tabular}
\caption{Basis comparison on five Church-numeral programs ($p$, $q$ free), run to normal
form. Each entry is mana with the \card{Flash Photography} engine / mana with the
\card{Replicator} engine G (MOVEs/COPYs, equal in both engines). Steps, SKI/BCKW/SKIBC:
11/10/5, 23/19/10, 22/16/10, 25/21/13, 33/25/15. Encoding sizes in Auras: 25/23/13,
49/41/23, 35/25/17, 27/23/15, 45/33/21. BCKW includes I; without it (I~=~$C\,K\,K$) BCKW
costs 92, 172, 236, 398 and 404 mana with the FP engine. Sources:
\texttt{sim/L6/output.txt}, \texttt{sim/L9/output.txt}.}\label{tab:basis}
\end{table}

\paragraph{Engines compared.} Over the 500 random terms the \card{Flash Photography} engine
spends 43{,}994 mana and 107{,}076 stack resolutions; the \card{Replicator} engine G spends
21{,}222 and 55{,}339, a 52\% saving in mana; engine R spends 144{,}396 and 368{,}368. On the
five programs, engine G saves 21--52\% on the SKI compilations and 8--35\% on the BCKW and
SKIBC compilations; the saving grows with the share of copies. Engine R costs 1.8 to 5.6
times the \card{Flash Photography} engine's mana on the programs. Physical scale is large:
$2\,2\,p\,q$ in SKI takes 634 mana, that is 634 Scepter loops and 1533 stack resolutions,
with the \card{Flash Photography} engine.

\section{What does not work}\label{sec:fails}

\subsection{The term in the library}

Our starting sketch put the term in the library in linear form, with the stack as
evaluation stack, phasing as suspension, token copies for W and milling for K. Six of its
nine claims fail. The library is hidden: ``Each library must be kept
in a single face-down pile. Players can't look at or change the order of cards in a
library'' (\CR{401.2}). Paper Magic has no effect that creates a card, and tokens cannot
leave the battlefield: ``A token that has left the battlefield can't move to another zone
or come back onto the battlefield'' (\CR{111.8}); copies of cards outside the stack and
battlefield cease to exist (\CR{707.10a}). Phasing applies only to permanents. We found no
trigger that reads a pattern of several library cards, and the tutors we checked shuffle
the library. What survives is a refereed, fuel-bounded two-stack machine: the term in
prefix order with an explicit application card, the library and graveyard as a zipper,
and face-up exile as a reserve of spare cards. A referee decides each move and counts
brackets. With a fixed deck it is a finite-state device. Its checker agreed with a
reference reducer on 14 hand-written and 2987 random normalising terms.

\subsection{The term on the stack}

The term can live on the stack as copies of inert instants in prefix code. \card{Reiterate}
copies an object to the top, and copies of a parked \card{Counterspell} with new targets
delete any object. This works cooperatively for all six combinators; its checker agreed
with a reference reducer on 10 fixed and 2986 random terms. But the term's structure lives
in the player: the extent of an argument is found by counting application markers, and no
card counts. C has no native stack analogue: no card we found reorders the stack, so a
reordering is a copy to the top plus a deletion, which makes C the most expensive rule.
A pointer design, with application nodes as spells whose targets are their subterms,
fails: new objects go only on top (\CR{405.2}), so the live spine must be re-copied above
every new node; a target that leaves the stack becomes illegal (\CR{608.2b}), so nothing
referenced can be reclaimed; and ``The same target can't be chosen multiple times for any
one instance of the word `target' '' (\CR{115.3}), so a node cannot point twice at one shared
subterm. That design was not simulated.

\subsection{Sharing}

An Aura has one host (\CR{303.4d}), so the attachment structure is a forest and subterms
cannot be shared. Copies carry no attachments (\CR{707.2}; rulings on \card{Flash
Photography}, \card{Replicator}, \card{Mirrormade} and \card{Copy Enchantment}). We found no
effect that copies a permanent together with its attachments; our searches were by
wording. So S and W copy node by node, and copies made for a duplicated argument may later
be thrown away by a K: $S\,I\,I\,(S\,I\,I)$ makes 96 copies in 40 steps.

\subsection{Forced reduction}

We asked how much of a reduction step forced play can do, in the sense of
\citet{cbh2019}: every step a state-based action or a mandatory trigger with uniquely
determined choices. With printed Oracle text, the answer we found is \emph{erasure} only:
the K cascade.

The obstruction has three assumptions. (E) The encoding is compositional: each subterm is a
subtree whose root is attached to the node encoding its context, and structure is carried
by attachment and card names only. (P) Printed text, no text-changing, and our card survey
of forced re-attachment effects is complete. (H) No addressing layer: no node carries a
position-dependent mark (control, colour, counters, tapped status, phasing) that would let a
forced choice single it out. Under these, the argument runs as follows. Pointers that
chain (an Aura on a non-creature permanent) are needed for depth at least 2, since an
attached permanent is never a creature (\CR{704.5p}, \CR{303.4d}). Every forced re-attacher
we found puts its object on a creature, a land or a player; the two routes onto a
non-creature (\card{Glamer Spinners}, and a token Aura entering under \CR{303.4f}) both
involve a choice. So below depth 1 forced play can only delete or suspend subtrees. Every
rule, even $I\,x \to x$, gives the surviving argument a new host. Hence, under (E), (P) and (H), no reduction
step at depth 2 or more can be forced. The steps that unbounded terms need chains and that every rule re-hosts the
surviving argument are arguments from (E) alone; the claim about re-attachers is a survey
result, and the survey was not exhaustive. At the root, one
forced K- or I-shaped step with real cards was sketched, using the fact that Equipment
survives its host's death; it does not compose. This sketch was not simulated.

Lifting (H) changes the picture in a model. A cursor machine in which every operation
touches only objects within one hop of a marked node (four event kinds: MARK, MOVE, COPY,
SAC; 14 cursors, not minimised) performs normal-order SKI and BCKW reduction; its checker
agreed with a reference reducer on 16 fixed and 3000 random terms. This is a model result.
A card realisation would need forced MOVE and COPY into non-creature hosts. The candidates
found are \card{Glamer Spinners}, made recurrent by \card{Followed Footsteps}, for MOVE, and
\card{Extravagant Replication} for COPY, with \card{Sterling Grove} to make targets unique.
Three conflicts are unresolved: \card{Sterling Grove}'s shroud would hide the very nodes
\card{Extravagant Replication} must target; \card{Glamer Spinners} moves all Auras of its
target, not one node; and the machine needs about 14 marks where control gives one bit.
Trigger ordering is a further hazard: with $K\,a\,\Omega$, one ordering choice decides
between termination and a mandatory loop, which is a draw (\CR{104.4b}).

\subsection{Phasing}

A phased-out permanent ``is treated as though it does not exist'' (\CR{702.26b}), and Auras
attached to a permanent that phases out phase out with it (\CR{702.26g}). A sweep of every
card whose Oracle text mentions phasing (76 cards, seven of them not legal in paper) found
no card that phases an Aura in a forced, repeatable way. Only four cards grant phasing,
and none reaches an Aura in paper. Chosen phase-out of any permanent exists (\card{Unite the
Coalition}, \card{The Pandorica}). Using it to suspend an argument only defers K's
clean-up: a cell of the discarded argument that was phased out directly phases in
unattached at its controller's next untap step, once its host is gone, and then it and its
subtree die (\CR{702.26i}, \CR{704.5m}). A suspended argument can be neither
moved nor copied. That a whole Aura-on-Aura subtree phases out together is implied by
\CR{702.26g} but not stated: the rule speaks of Auras attached to ``that permanent'', and
gives no example with an Aura on an Aura. The sweep is exhaustive only with respect to the
wordings searched.

\section{Threats to validity and open rules questions}\label{sec:threats}

\begin{itemize}[nosep,leftmargin=*]
\item No judge or official source has confirmed any of the rules readings. All legality
  claims are on our reading of the Comprehensive Rules of 6 November 2026 and the cited
  rulings.
\item That attachments ride along when an Aura is moved is an argument from absence
  (\CR{701.3}, no zone change, \CR{400.7}), not a stated rule.
\item Whether attachment cycles are legal is unresolved. \CR{303.4d} forbids only
  self-attachment, so a MOVE could close a cycle on our reading. The protocol never creates
  one, and the simulators reject such a move only as a protocol error.
\item The \card{Replicator} seeding step rests on a reading of two rulings, and the card is
  new (released 13 November 2026).
\item The setup of the initial board and of the engine is assumed, not constructed.
\item The opponent is assumed to pass. Every resolution is a priority round at which the
  opponent could respond.
\item Card surveys (copy effects, Aura movers, forced re-attachers, phasing, labels) were by
  Oracle wording on Scryfall. A card phrased differently could exist.
\item The simulators check the abstract moves and some rule mechanics. They do not check the
  same-controller clause, cycle legality or trigger timing.
\item The cooperative model does not stop a player from making other legal moves, for
  example moving a node onto its own descendant, onto \card{Auratog} (which \card{Bound in Gold}
  would then switch off), or sacrificing a cell mid-term. Such moves break the computation,
  not the rules.
\end{itemize}

\section{Conclusion and open problems}

Combinatory-logic terms can be laid out on the battlefield as forests of Auras, and
normal-order reduction for BCKW and for SKI becomes three kinds of ordinary game action.
B and C rearrange attachments, W and S copy, and K hands its discarded argument to the
game's own state-based actions. The construction uses a small set of real paper cards,
and the board can be read as the term at every step. All of this holds in a cooperative
model and on our reading of the rules.

Open problems: (1) \emph{Forced reduction.} Can MARK, MOVE and COPY into non-creature hosts
be realised by mandatory triggers with unique choices, and at what cost in cards? (2) \emph{A
mechanised model.} A formal model of the relevant rules (attachment, copying, state-based
actions), in which the proposition of Section~\ref{sec:machine} would be a theorem about
the game rather than about the abstract moves. (3) \emph{Sharing.} An effect that copies
a permanent with its attachments, or any native relation with out-degree above one, would
allow graph reduction; we found none.

\bibliographystyle{plainnat}
\bibliography{paper}

\appendix

\section{Cards}\label{app:cards}

Oracle text as served by Scryfall on 21 November 2026. A source check confirmed the texts
of the first 24 cards below; \card{Lithoform Engine} is quoted from a separate Scryfall fetch
that the check did not cover. ``(\dots)'' marks reminder text omitted in our notes. The 24
checked cards are paper cards legal in Legacy, Vintage and Commander according to Scryfall;
\card{Lithoform Engine} is a paper card legal in Legacy according to the same source.

\begingroup\small
\begin{description}[leftmargin=1em,itemsep=1pt,parsep=0pt]
\item[\card{Island}] Basic Land. ``(\ms{T}: Add \ms{U}.)''
\item[\card{Simic Guildmage}] \ms{G/U}\ms{G/U}, Creature. ``(\ms{G/U} can be paid with either
  \ms{G} or \ms{U}.) / \ms{1}\ms{G}: Move a +1/+1 counter from target creature onto another
  target creature with the same controller. / \ms{1}\ms{U}: Attach target Aura attached to a
  permanent to another permanent with the same controller.''
\item[\card{Auratog}] \ms{1}\ms{W}, Creature. ``Sacrifice an enchantment: This creature gets
  +2/+2 until end of turn.''
\item[\card{Flash Photography}] \ms{2}\ms{U}\ms{U}, Sorcery. ``You may cast this spell as though
  it had flash if it targets a permanent you control. / Create a token that's a copy of
  target permanent. / Flashback \ms{4}\ms{U}\ms{U}''
\item[\card{Reiterate}] \ms{1}\ms{R}\ms{R}, Instant. ``Buyback \ms{3} (You may pay an
  additional \ms{3} as you cast this spell. If you do, put this card into your hand as it
  resolves.) / Copy target instant or sorcery spell. You may choose new targets for the
  copy.''
\item[\card{Isochron Scepter}] \ms{2}, Artifact. ``Imprint --- When this artifact enters, you
  may exile an instant card with mana value 2 or less from your hand. / \ms{2}, \ms{T}: You
  may copy the exiled card. If you do, you may cast the copy without paying its mana
  cost.''
\item[\card{Dramatic Reversal}] \ms{1}\ms{U}, Instant. ``Untap all nonland permanents you
  control.''
\item[\card{Gilded Lotus}] \ms{5}, Artifact. ``\ms{T}: Add three mana of any one color.''
\item[\card{Replicator}] \ms{2}, Artifact. ``When this artifact enters, create a Food token.
  (It's an artifact with `\ms{2}, \ms{T}, Sacrifice this token: You gain 3 life.') /
  \ms{2}, \ms{T}: Create a tapped token that's a copy of target token you control.''
\item[\card{Mythos of Illuna}] Sorcery. ``Create a token that's a copy of target permanent.''
  Quoted in part; an \ms{R}\ms{G} fight clause is omitted.
  % TODO(author): full Oracle text of Mythos of Illuna is not in the notes.
\item[\card{Replication Technique}] Sorcery. ``Create a token that's a copy of target permanent
  you control.'' Quoted in part; the demonstrate line is omitted.
  % TODO(author): full Oracle text of Replication Technique is not in the notes.
\item[\card{Indestructibility}] \ms{3}\ms{W}, Aura. ``Enchant permanent / Enchanted permanent
  has indestructible. (Effects that say ``destroy'' don't destroy that permanent. A
  creature with indestructible can't be destroyed by damage.)''
\item[\card{Bound in Gold}] \ms{2}\ms{W}, Aura. ``Enchant permanent / Enchanted permanent can't
  attack, block, or crew Vehicles, and its activated abilities can't be activated unless
  they're mana abilities.''
\item[\card{Intercessor's Arrest}] \ms{2}\ms{W}, Aura. ``Enchant permanent / Enchanted permanent
  can't attack, block, or crew Vehicles. Its activated abilities can't be activated unless
  they're mana abilities.''
\item[\card{Curator's Ward}] \ms{2}\ms{U}, Aura. ``Enchant permanent / Enchanted permanent has
  hexproof. / When enchanted permanent leaves the battlefield, if it was historic, draw two
  cards. (Artifacts, legendaries, and Sagas are historic.)''
\item[\card{Essence Leak}] \ms{U}, Aura. ``Enchant permanent / As long as enchanted permanent is
  red or green, it has `At the beginning of your upkeep, sacrifice this permanent unless you
  pay its mana cost.' ''
\item[\card{Faith's Fetters}] \ms{3}\ms{W}, Aura. ``Enchant permanent / When this Aura enters,
  you gain 4 life. / Enchanted permanent can't attack or block, and its activated abilities
  can't be activated unless they're mana abilities.''
\item[\card{Elemental Resonance}] \ms{2}\ms{G}\ms{G}, Aura. ``Enchant permanent / At the
  beginning of your first main phase, add mana equal to enchanted permanent's mana cost.
  (Mana cost includes color. If a mana symbol has multiple colors, choose one.)''
\item[\card{Confiscate}] \ms{4}\ms{U}\ms{U}, Aura. ``Enchant permanent / You control enchanted
  permanent.''
\item[\card{Lay Claim}] \ms{5}\ms{U}\ms{U}, Aura. ``Enchant permanent / You control enchanted
  permanent. / Cycling \ms{2} (\dots)''
  % TODO(author): cycling reminder text of Lay Claim is elided in the notes.
\item[\card{Volition Reins}] \ms{3}\ms{U}\ms{U}\ms{U}, Aura. ``Enchant permanent / When this Aura
  enters, if enchanted permanent is tapped, untap it. / You control enchanted permanent.''
\item[\card{Take Possession}] \ms{5}\ms{U}\ms{U}, Aura. ``Split second (As long as this spell is
  on the stack, players can't cast spells or activate abilities that aren't mana
  abilities.) / Enchant permanent / You control enchanted permanent.''
\item[\card{Psychic Overload}] \ms{3}\ms{U}, Aura. ``Enchant permanent / When this Aura enters,
  tap enchanted permanent. / Enchanted permanent doesn't untap during its controller's untap
  step. / Enchanted permanent has `Discard two artifact cards: Untap this permanent.' ''
\item[\card{Dream Leash}] \ms{3}\ms{U}\ms{U}, Aura. ``Enchant permanent / You can't choose an
  untapped permanent as this spell's target as you cast it. / You control enchanted
  permanent.''
\item[\card{Lithoform Engine}] \ms{4}, Legendary Artifact. ``\ms{2}, \ms{T}: Copy target
  activated or triggered ability you control. You may choose new targets for the copy. /
  \ms{3}, \ms{T}: Copy target instant or sorcery spell you control. You may choose new targets
  for the copy. / \ms{4}, \ms{T}: Copy target permanent spell you control. (The copy becomes a
  token.)''
\end{description}
\endgroup

\section{Rules quoted}\label{app:rules}

Comprehensive Rules, edition effective 6 November 2026. Ellipses mark omissions.

\begingroup\small
\begin{description}[leftmargin=1em,itemsep=1pt,parsep=0pt]
\item[104.4b] ``\dots\ Loops that contain an optional action don't result in a draw.''
\item[106.4] ``Each player's mana pool empties at the end of each step and phase''.
\item[110.1] ``A permanent is a card or token on the battlefield.''
\item[111.2] ``The player who creates a token is its owner. The token enters the battlefield
  under that player's control.''
\item[111.3] ``The spell or ability that creates a token may define the values of any number
  of characteristics for the token. \dots\ The characteristic values defined this way are
  functionally equivalent to the characteristic values that are printed on a card; for
  example, they define the token's copiable values.''
\item[111.7] ``A token that's in a zone other than the battlefield ceases to exist. This is a
  state-based action'' \dots
\item[111.8] ``A token that has left the battlefield can't move to another zone or come back
  onto the battlefield. If such a token would change zones, it remains in its current zone
  instead. It ceases to exist the next time state-based actions are checked; see rule
  704.''
\item[115.3] ``The same target can't be chosen multiple times for any one instance of the
  word `target' on a spell or ability.'' \dots
\item[117.5] ``Each time a player would get priority, the game first performs all applicable
  state-based actions as a single event (see rule 704, `State-Based Actions'), then repeats
  this process until no state-based actions are performed.'' \dots
\item[303.4d] ``An Aura can't enchant itself. \dots\ An Aura can't enchant more than one object
  or player.''
\item[303.4f] ``If an Aura is entering the battlefield under a player's control by any means
  other than by resolving as an Aura spell, and the effect putting it onto the battlefield
  doesn't specify the object or player the Aura will enchant, that player chooses what it
  will enchant as the Aura enters the battlefield. The player must choose a legal object or
  player according to the Aura's enchant ability and any other applicable effects.''
\item[400.7] ``An object that moves from one zone to another becomes a new object with no
  memory of, or relation to, its previous existence.''
\item[401.2] ``Each library must be kept in a single face-down pile. Players can't look at or
  change the order of cards in a library.''
\item[405.2] ``Each time an object is put on the stack, it's put on top of all objects
  already there.''
\item[405.5] ``If the stack is empty when all players pass, the current step or phase ends
  and the next begins.''
\item[608.2b] ``If the spell or ability specifies targets, it checks whether the targets are
  still legal. A target that's no longer in the zone it was in when it was targeted is
  illegal. \dots\ If all its targets, for every instance of the word `target,' are now
  illegal, the spell or ability doesn't resolve.'' \dots
\item[701.3a] ``To attach an Aura, Equipment, or Fortification to an object or player means
  to take it from where it currently is and put it onto that object or player. \dots\ An
  Aura, Equipment, or Fortification can't be attached to an object or player it couldn't
  enchant, equip, or fortify, respectively.''
\item[701.3c] ``Attaching an Aura, Equipment, or Fortification on the battlefield to a
  different object or player causes the Aura, Equipment, or Fortification to receive a new
  timestamp.''
\item[701.21a] ``To sacrifice a permanent, its controller moves it from the battlefield
  directly to its owner's graveyard. A player can't sacrifice something that isn't a
  permanent, or something that's a permanent they don't control. Sacrificing a permanent
  doesn't destroy it''.
\item[702.26b] ``If a permanent phases out, its status changes to `phased out.' Except for
  rules and effects that specifically mention phased-out permanents, a phased-out permanent
  is treated as though it does not exist. It can't affect or be affected by anything else
  in the game.'' \dots
\item[702.26g] ``When a permanent phases out, any Auras, Equipment, or Fortifications attached
  to that permanent phase out at the same time. This alternate way of phasing out is known
  as phasing out `indirectly.' An Aura, Equipment, or Fortification that phased out
  indirectly won't phase in by itself, but instead phases in along with the permanent it's
  attached to.''
\item[702.26i] ``An Aura, Equipment, or Fortification that phased out directly will phase in
  attached to the object or player it was attached to when it phased out, if that object is
  still in the same zone or that player is still in the game. If not, that Aura, Equipment,
  or Fortification phases in unattached. State-based actions apply as appropriate.''
\item[702.27a] (buyback) ``put this spell into its owner's hand instead of into that player's
  graveyard as it resolves.''
\item[704.3] ``Whenever a player would get priority \dots, the game checks for any of the
  listed conditions for state-based actions, then performs all applicable state-based
  actions simultaneously as a single event. If any state-based actions are performed as a
  result of a check, the check is repeated'' \dots
\item[704.5m] ``If an Aura is attached to an illegal object or player, or is not attached to
  an object or player, that Aura is put into its owner's graveyard.''
\item[707.2] ``The copiable values are the values derived from the text printed on the
  object \dots\ Other effects (including type-changing and text-changing effects), status,
  counters, and stickers are not copied.''
\item[707.10] ``To copy a spell, activated ability, or triggered ability means to put a copy
  of it onto the stack; a copy of a spell isn't cast and a copy of an activated ability
  isn't activated. A copy of a spell or ability copies both the characteristics of the
  spell or ability and all decisions made for it, including modes, targets, the value of X,
  and additional or alternative costs. \dots\ A copy of a spell or ability is controlled by
  the player under whose control it was put on the stack. A copy of a spell is itself a
  spell, even though it has no spell card associated with it.''
\item[707.10a] ``If a copy of a spell is in a zone other than the stack, it ceases to exist.
  If a copy of a card is in any zone other than the stack or the battlefield, it ceases to
  exist.'' \dots
\item[707.10c] ``\dots\ If the player chooses to change some or all of the targets, the new
  targets must be legal. Once the player has decided what the copy's targets will be, the
  copy is put onto the stack with those targets.''
\item[707.10f] ``Some effects copy a permanent spell. As that copy resolves, it ceases being a
  copy of a spell and becomes a token permanent.'' \dots
\end{description}
\endgroup

\end{document}
